\documentclass{atscv}
\usepackage{atscv-schema-support}
\usepackage{atscv-color-burgundy}
\usepackage{atscv-style-consulting}
\usepackage{atscv-layout-letter}
\usepackage{atscv-lang-de}
\begin{document}
\cvname{Fabio Schmeil}
\cvtitle{Bewerbung als Customer Support Specialist DACH}
\cvemail{f.schmeil@gmail.com}
\cvphone{+49 176 57624160}
\cvlocation{Deutschland | Remote | Umzugsbereit}
\makecvheader

\cvsection{Anschreiben}
Sehr geehrtes Hiring-Team,

mit großem Interesse bewerbe ich mich für die Position als Customer Support Specialist DACH im Umfeld von Wolters Kluwer, Libra und Legal-AI-nahen Produkten. Ich arbeite gerne strukturiert an Supportfällen, priorisiere sauber nach Auswirkung und Dringlichkeit und sichere eine nachvollziehbare Kommunikation über den gesamten Bearbeitungszyklus.

Im operativen Alltag verbinde ich Ticket- und Eskalationsbearbeitung mit konsequenter CRM-Pflege, präziser Wissensdatenbank-Dokumentation und enger Abstimmung mit Customer Success sowie Produktteams. So werden wiederkehrende Fehlerbilder nicht nur gelöst, sondern in belastbares Produktfeedback übersetzt.

Besonders wichtig ist mir eine systematische Ursachenanalyse: Muster erkennen, Hypothesen prüfen, Maßnahmen bewerten und Ergebnisse transparent machen. Meine Arbeitsweise lässt sich als ausgeprägtes Mustererkennen, hohe Detailgenauigkeit und schnelle Strukturierung komplexer Zusammenhänge zusammenfassen. Bei Bedarf nutze ich zudem systemische Analyse von Verhaltens- und Risikomustern in komplexen technischen Umgebungen, um Eskalationen früh zu entschärfen.

SaaS- und AI-Produkte motivieren mich besonders, weil sie kontinuierliches Lernen mit direktem Kundennutzen verbinden. Genau hier möchte ich meine Support-Erfahrung, meine klare schriftliche Kommunikation und meine Zusammenarbeit mit Legal Engineering, Produkt und Operations einbringen.

Ich freue mich auf ein Gespräch.

Mit freundlichen Grüßen

Fabio Schmeil
\end{document}
